% =====================================================================
%  Introduction — revised against the outline
%
%  WHAT CHANGED (everything else is the previous text, largely verbatim):
%   1. The three prior-work paragraphs (scGen/ContrastiveVI; SAMS-VAE;
%      SVAE/SVAE+) are merged into two, grouped by what the methods do
%      rather than one paragraph per baseline.
%   2. The gap paragraph now lists four requirements rather than three,
%      adding combinatorial interactions so it matches the four
%      dimensions introduced two paragraphs earlier and the outline.
%   3. The SHERLOCK paragraph describes the additive displacement of the
%      background used by the model (see Methods). The background is also
%      no longer described as condition-dependent.
%
%  Text that differs from the previous version is wrapped in { }. These
%  are LaTeX groups, so they are visible in the source and invisible in
%  the output; strip them once the changes are approved.
%
%  STILL TO RESOLVE: the [CITE ...] placeholders are carried over
%  unchanged. Elham's outline block is retained at the top for comparison
%  and can be deleted once the text is approved.
%
%  NOTE: the sentence in Results 2.1 describing the latent transition
%  still needs replacing to match this paragraph.
% =====================================================================

\section{Introduction}
\label{sec:intro}
\ea{outline: Para 1: The opportunity: Single-cell perturbation technologies (CRISPR, chemical, and spatial perturbation screens) now enable systematic interrogation of gene function at unprecedented scale, creating an urgent need for computational methods that extract mechanistic insight from these large interventional datasets.
Para2 – The bio questions: A single perturbation screen should ideally answer four fundamental questions: (i) which perturbations share mechanisms, (ii) what downstream genes and pathways they affect, (iii) how their effects change across experimental conditions, and (iv) how perturbations interact in combinations.
Para3 – The computational gap: Existing methods address these questions separately—focusing on perturbation prediction, differential expression, representation learning, or combinatorial perturbations—but no single framework enables all of these analyses from a common perturbation representation.
Para4 – The gap: Achieving this requires learning structured representations of perturbations that capture shared biological mechanisms while simultaneously supporting counterfactual estimation of downstream perturbation effects, quantification of condition-dependent responses, and analysis of combinatorial interactions.
Para5 –  solution: An interpretable generative framework that learns structured perturbation representations to organize perturbations into mechanistic groups, infer counterfactual perturbation-to-gene effect networks, quantify condition-dependent perturbation responses, and model combinatorial perturbation effects across diverse single-cell perturbation modalities.}

Large-scale perturbation assays, including CRISPR-based, chemical, and spatial screens, enable
systematic characterization of cellular responses and gene function. These experiments generate
high-dimensional measurements of transcriptional changes induced by diverse perturbations across
heterogeneous biological and experimental contexts. The increasing scale and complexity of such
datasets require computational methods that can move beyond simply predicting post-perturbation gene
expression in and out-of-distributions and infer perturbation-response relationships to extract
causal mechanistic structure from interventional data.

An ideal perturbation screen should support systematic characterization of molecular systems along
four complementary dimensions. First, it should quantify the downstream effects of each perturbation,
including changes in gene expression, signaling pathways, regulatory programs, and cellular
phenotypes, while resolving variation across cell types, states, and experimental contexts. Second,
it should identify similarities and differences among perturbations by comparing their response
profiles and determining whether they act through shared, partially overlapping, or distinct
mechanisms of action. Third, it should characterize how exogenous variables, including environmental
conditions, treatment settings, temporal context, and other measured covariates, modify the
magnitude, direction, and specificity of perturbation effects. Fourth, it should model interactions
among multiple perturbations under combinatorial interventions, including additive, synergistic,
antagonistic, and context-dependent effects. Joint analysis of these four dimensions is required to
infer perturbation-response mechanisms, prioritize candidate therapeutic targets, identify
context-specific vulnerabilities, and characterize the structure and dynamics of cellular state
transitions.

Existing computational approaches for single-cell perturbation analysis typically address these
objectives in isolation and often do not explicitly represent the mechanistic structure required for
causal effect estimation under interventions on exogenous variables. In addition, increasing reliance
on high capacity models such as foundation models and large language models[CITE GEAR, scGPT,
cell2sentence, etc.], can limit interpretability of perturbation-response mechanisms and introduce
statistical and computational challenges related to identifiability, over-parameterization, and
overfitting. These limitations are particularly consequential for counterfactual inference, which
requires the model to first infer the latent exogenous noise variables that are compatible with the
observed data, then apply alternative do-interventions while holding the abducted noise realization
fixed, and finally propagate the intervened system through the structural equations to estimate
counterfactual outcomes. Without an identifiable latent-variable parameterization and a separation
between exogenous variation and intervention-specific mechanisms, predictive models of
post-perturbation expression cannot, in general, support counterfactual estimation at the higher
levels of Pearl's causal hierarchy[CITE Elias and Pearl, 2020 paper and Pearl 1999 paper].

Even if we zoom in to more interpretable frameworks such as deep generative models, existing
VAE-based approaches do not jointly address these objectives. scGen\cite{lotfollahi2019scgen}
represents a perturbation as a global displacement vector estimated by latent-space arithmetic
between control and perturbed populations{, a formulation that assumes} the perturbation effect can be
summarized by a context-invariant translation and {that} generally models each perturbation
independently. ContrastiveVI\cite{weinberger2023contrastivevi} {instead }decomposes cellular variation
into shared and salient latent components to distinguish background variation from
condition-associated variation{; this binary decomposition, however,} does not assign individual
perturbations to specific latent mechanisms or characterize which transcriptional programs, pathways,
or regulatory modules mediate their effects. {Neither identifies perturbation-specific regulatory
programs, represents dependencies among perturbations, or specifies how exogenous variables modify
the perturbation mechanism, and neither provides a structural representation of perturbation
interactions.}

{Sparse mechanism-shift and structured latent-variable models come closer.}
SAMS-VAE\cite{bereket2023samsvae} introduces perturbation-specific masks over latent variables,
allowing each intervention to modify a subset of latent dimensions{, while
SVAE\cite{lachapelle2022mechanism} and SVAE+\cite{lopez2023sviaeplus} impose additional structure or
supervision on the latent representation to improve correspondence between latent variables and
generative factors}. Although {these parameterizations} encode modularity and can improve sparsity of
the inferred perturbation effects, the latent variables remain non-identifiable in general: multiple
latent parameterizations may induce the same observed-data distribution while assigning different
biological interpretations to the masked dimensions. Sparsity {or supervision }alone therefore does not
guarantee recovery of the underlying causal mechanisms, regulatory programs, or intervention targets.
Moreover, perturbation embeddings are typically {parameterized independently, without} an explicit
relational structure linking perturbations through shared targets, pathways, or mechanisms{, so these
models} do not directly represent mechanistic similarity among interventions, context-dependent
modulation of perturbation effects, or higher-order interactions between perturbations. More
generally, without an identifiable structural causal model that separates exogenous variation,
endogenous cellular state, and intervention-specific mechanisms, these VAE formulations cannot
implement counterfactual inference through abduction, intervention, and prediction while holding the
inferred exogenous noise realization fixed.

Most fundamentally, existing methods struggle with interpretability: the learned latent dimensions
often lack clear biological meaning, making it difficult to translate computational findings into
testable mechanistic hypotheses. Methods that enforce sparsity frequently do so at the expense of
predictive accuracy, whereas methods optimized for prediction tend to produce entangled
representations in which perturbation effects are diffusely distributed across latent dimensions. As
a result, the field still lacks a unified framework that can simultaneously: (1) learn structured
groupings of perturbations that reflect underlying mechanistic relationships; (2) robustly link
genetic perturbations to their downstream transcriptional consequences to enable mechanistic insight;
(3) explicitly model condition-dependent perturbation effects{; and (4) represent combinatorial
interventions in a way that exposes how perturbations interact}. This framework should generalize
across diverse perturbation settings, including different CRISPR modalities (CRISPRi, CRISPRa, and
CRISPR knockout), as well as across single-cell and spatial data modalities, while supporting
biological discovery and hypothesis generation.

Here, we introduce SHERLOCK, a variational autoencoder framework that addresses these limitations. We
represent perturbations as structured interventions on a shared latent cellular state, in order to
have both a generative view and a structural causal model(SCM) view of our model to allow
counterfactual estimation. Each perturbation is parameterized by a sparse latent shift,
\[
\Delta_p = W_p \odot A_p,
\]
where the gate selects affected latent coordinates and the low-rank covariance over $A$ captures
dependencies among perturbation mechanisms. {A background variable ($z_0$), inferred from control
cells, is displaced additively by ($\Delta_p$)}, separating contextual variation from
intervention-specific effects. This structure also allows explicit abduction step so we can infer
joint distribution of exogenous noise and hold them fixed. A shared negative-binomial decoder maps
intervened latent states to gene-expression distributions, enabling perturbation-to-gene effects to
be defined as contrasts between decoded potential outcomes. The model extends to combinatorial
interventions through additive perturbation shifts and a learned interaction residual. We specify
causal estimands induced by the SCM view of our model and estimates them under stated structural and
statistical assumptions, and prove theoretical results such as causal estimand identifiability while
showing empirical evidence of reproducibility.

Applied to the genome-wide Perturb-seq dataset generated by Replogle et al.,
\cite{replogle2022mapping}, our approach revealed coherent grouping of perturbations according to
pathway annotations, demonstrating that it captures ground truth biological grouping and pathway
structures more faithfully than existing methods, achieves interpretable separation of underlying
mechanisms, and generalizes to diverse experimental conditions. In addition, it demonstrated strong
performance across multiple perturbation modalities, including drug treatment\cite{giglio2025heterogeneous},
CRISPR inhibition \cite{replogle2022mapping}, activation\cite{shi2026drug}, and RAEFISH spatial
perturbation datasets\cite{cell2025spatial}. Overall, the framework supports meaningful grouping of
perturbations, the construction of perturbation-to-target gene graphs, {identification of conditional
perturbation response, and quantification of genetic interactions between combined perturbations},
providing a robust platform for both exploratory analysis and principled hypothesis testing in
perturbation biology. Through our approach, we aim to bridge data-driven modeling with mechanistic
understanding, advancing the field's capacity for biological discovery in high-dimensional
experimental data.
